\section{Third Order Orthogonality}\label{sec: Third Order Orth}

In this section, we assume 
\begin{equation}\label{eq: arity 6 assumption}
    f\in \mathcal{F} \text{ is arity } 6, \text{irreducible, satisfies {\sc 2nd-Orth} and } B(f) \text{ closure \cref{prop: binary closure}.}
\end{equation}

We are going to prove that $f$ satisfies {\sc 3rd-Orth}, i.e., $\ket{f}$ is an AME$(6,2)$ state.
The proof enumerates all possible cases of the finite group $B(f)\le \mathbf{SO}(3)$.
We fix arbitrary distinct indices $\{i,j,k\}$ and show that either $\rho_{ijk}=\frac{1}{8}I$ or \eqref{eq: arity 6 assumption} cannot be satisfied.

\subsection{Polyhedral Groups}\label{subsec: polyhedral groups}

In this subsection, we consider the case that $B(f)\le \mathbf{SO}(3)$ is a polyhedral group.
It can be the tetrahedral group $A_4$, the octahedral group $S_4$, or the icosahedral group $A_5$.

Let $\rho=\ket{f}\bra{f}$ be the density matrix of $f$.
By merging we can realize the state with marginal density matrix in Pauli representation:
$$\rho_{ijk}=\frac{1}{8}\sum_{p,q,r=0}^3 t_{pqr} \sigma_p\otimes\sigma_q\otimes\sigma_r.$$
By {\sc 2nd-Orth}, we have 
\begin{lemma}\label{lm: Pauli rep by 2nd-Orth}
    $$\rho_{ijk}=\frac{1}{8}\left(I^{\otimes3}+\sum_{p,q,r=1}^3 t_{pqr} \sigma_p\otimes\sigma_q\otimes\sigma_r\right).$$
\end{lemma}
\begin{proof}
    First, $t_{000}=\mathrm{Tr}(\rho_{ijk}(\sigma_0\otimes\sigma_0\otimes\sigma_0))=\mathrm{Tr}(\rho_{ijk})=1.$
    We only need to prove $t_{pqr}=0$ if there are $1$ or $2$ zeros among $\{p,q,r\}.$
    For example,
    $t_{120}=\mathrm{Tr}(\rho_{ijk}(\sigma_1\otimes \sigma_2\otimes I))=\mathrm{Tr}(\rho_{ij}(\sigma_1\otimes \sigma_2)).$
    By {\sc 2nd-Orth}, $\rho_{ij}=\frac{1}{4}I$.
    So $t_{120}=\frac{1}{4}\mathrm{Tr}(\sigma_1\otimes \sigma_2)=\frac{1}{4}\mathrm{Tr}(\sigma_1)\mathrm{Tr}(\sigma_2)=0.$
    Similarly, $t_{pqr}=0$ if there are $1$ or $2$ zeros among $\{p,q,r\}.$
\end{proof}


So $f$ satisfies $3$rd-Orth iff $t_{pqr}=0$ for all $p,q,r\in\{1,2,3\}.$

\begin{definition}
    For every $b\in B(f)$ with matrix $M(b)\in \mathbf{SU}(2)$, 
    define $\operatorname{vec}(b)=(b_{00},b_{01},b_{10},b_{11})^{\tt T}$, $R_b\in \mathbf{SO}(3)$ be the adjoint rotation of $b$, and the matrix $C(b)\in \mathbb{R}^{3\times 3}$ by $C(b)_{pq}=\operatorname{vec}(b)^\dagger (\sigma_p\otimes\sigma_q) \operatorname{vec}(b)$ for $p,q\in\{1,2,3\}$.
\end{definition}

\begin{lemma}\label{lm: Cb and Rb}
    $C(b)=2R_{b^{-1}}^{\tt T}S$, where $S=\mathrm{diag}(1,-1,1).$
\end{lemma}
\begin{proof}
    %Suppose the binary signature $b$ is normalized with matrix $M(b)\in \mathbf{SU}(2).$
    Recall the correspondence between $M(b)\in\mathbf{SU}(2)$ and $R_b\in \mathbf{SO}(3)$ in \cref{para: SU2 and SO3} gives
    $$M(b) \sigma_pM(b)^\dagger=\sum_{r=1}^3(R_b)_{rp}\sigma_r\implies M(b)^\dagger\sigma_pM(b)=\sum_{r=1}^3(R_{b^{-1}})_{rp}\sigma_r$$ 
    for $p=1,2,3.$
    Let $s_1=1, s_2=-1, s_3=1$, notice that $s_q\sigma_q=\sigma_q^{\tt T}$ for $q=1,2,3.$
    Then
    \begin{equation*}
        \begin{aligned}
            C(b)_{pq}=&\operatorname{vec}(b)^\dagger(\sigma_p\otimes\sigma_q) \operatorname{vec}(b)=\mathrm{Tr}(M(b)^\dagger \sigma_p M(b)\sigma_q^{\tt T})
            =\mathrm{Tr}\left(\sum_{r=1}^3(R_{b^{-1}})_{rp}\sigma_rs_q\sigma_q\right)
            =s_q \sum_{r=1}^3 (R_{b^{-1}})_{rp}\mathrm{Tr}(\sigma_r\sigma_q)\\
            =&2s_q \sum_{r=1}^3 (R_{b^{-1}})_{rp}\delta_{rq}
            =2s_q(R_{b^{-1}})_{qp}
            =(2R_{b^{-1}}^{\tt T}S)_{pq}.
        \end{aligned}
    \end{equation*}
\end{proof}

\begin{lemma}[Pauli slice]\label{lm: Pauli slice}
    For every $b\in B(f)$ and $r=1,2,3$, we have $\sum_{p,q=1}^3C(b)_{pq}t_{p,q,r}=0.$
\end{lemma}

\begin{proof}
    Suppose after merging $f$ and $\overline{f}$, the dangling edges are $x_i,x_j,x_k$ and $y_i,y_j,y_k$.
    Since $b\in B(f)$, we have $\overline{b}\in B(f)$.
    Connect $x_i$ and $x_j$ using $\overline{b}$, $y_i$ and $y_j$ using $b$, we get a binary signature $g(x_k,y_k)\in \mathcal{U}$ by binary closure \cref{prop: binary closure}.
    We have
    \begin{equation*}
        \begin{aligned}
            \mu I_2
            =&M(g)M(g)^{\dagger}=(\operatorname{vec}(b)^\dagger \otimes I_2) \rho_{ijk}(\operatorname{vec}(b) \otimes I_2)\\
            =&\frac{1}{8}(\operatorname{vec}(b)^\dagger \otimes I_2)\left(I_2^{\otimes3}+\sum_{p,q,r=1}^3 t_{pqr} \sigma_p\otimes\sigma_q\otimes\sigma_r\right)(\operatorname{vec}(b) \otimes I_2)\\
            =& \frac{1}{8}\left(\operatorname{vec}(b)^{\dagger}\operatorname{vec}(b)\otimes I_2+\sum_{p,q,r=1}^3t_{pqr}(\operatorname{vec}(b)^\dagger (\sigma_p\otimes\sigma_q) \operatorname{vec}(b))\sigma_r\right)\\
            =&\frac{|\operatorname{vec}(b)|^2}{8}I_2+\frac{1}{8}\sum_{r=1}^3\left(\sum_{p,q=1}^3 C(b)_{pq}t_{pqr}\right)\sigma_r.&
        \end{aligned}
    \end{equation*}
    Since $\{I_2,\sigma_1,\sigma_2,\sigma_3\}$ is a basis of $\mathbb{C}^{2\times 2}$, we have $\sum_{p,q=1}^3 C(b)_{pq}t_{pqr}=0$ for $r=1,2,3.$
\end{proof}
\begin{remark}
    For $r=1,2,3$, we define $T^{(r)}$ be the $3\times 3$ matrix with entries $T^{(r)}_{pq}=t_{pqr}$.
    This is a slice of the the Pauli coefficients $t_{pqr}.$
    \cref{lm: Pauli slice} says $T^{(r)}$ is orthogonal to $C(b)$ under the Frobenius inner product in $M_3(\mathbb{R}).$
    Similarly, we can define $T^{(p)}$ and $T^{(q)}$, and they are also orthogonal to $C(b)$.
\end{remark}

The following lemma utilizes the assumption that $B(f)=A_4,S_4$ or $A_5$.
\begin{lemma}\label{lm: rep lemma}
    Let $W=\mathrm{Span}_{\mathbb{R}}\{C(b)\mid b\in B(f)\}$, then $W=M_3(\mathbb{R})$.
\end{lemma}
\begin{proof}
    By Section 4.8 and Theorem 4.5.1 in \cite{reptheory}, for $B(f)=A_4,S_4,A_5$, the natural three-dimensional rotation representation $\rho:b\mapsto R_b$ is irreducible.
    Apply the density theorem (Theorem 3.2.2 in \cite{reptheory}), we have $\mathrm{Span}_{\mathbb{C}}\{R_b\mid b\in B(f)\}=M_3(\mathbb{C})$.
    Since all matrices $R_b$ are real, $\mathrm{Span}_{\mathbb{R}}\{R_b\mid b\in B(f)\}=M_3(\mathbb{R})$.
    By \cref{lm: Cb and Rb}, $W=M_3(\mathbb{R}).$
\end{proof}

\begin{lemma}\label{lm: arity 6 polyhedral group 3rd orth}
    %$t_{p,q,r}=0$ for all $p,q,r\in\{1,2,3\}$. i.e., $f$ satisfies {\sc 3rd-Orth}.
    If $B(f)$ is isomorphic to a polyhedral group, then $f$ satisfies {\sc 3rd-Orth}.
\end{lemma}
\begin{proof}
    By \cref{lm: Pauli slice}, the matrix $T^{(r)}$ with entries $(T^{(r)})_{pq}=t_{pqr}$ is in $W^{\perp}$. 
    Here we use the Frobenius inner product in $M_3(\mathbb{R})$: $\braket{A,B}=\mathrm{Tr}(A^TB).$
    By \cref{lm: rep lemma}, $W^\perp=\{O\}.$ 
    Thus, $T^{(r)}=O$ for every $r=1,2,3$, or $t_{pqr}=0$ for every $p,q,r\in\{1,2,3\}.$
\end{proof}

\subsection{Dihedral Groups of Order at Least Six}
In this subsection, we assume that $B(f)\cong D_n(n\ge 3)$ which is the dihedral group of order $2n.$
We still use the Pauli representation of $\rho_{ijk}$ in \cref{subsec: polyhedral groups}.

Suppose in $\mathbb{R}^3$, the regular $n$-polygon $P_n$ lies on the $x$-$y$ plane, and $P_n$ is symmetric relative to the $x$-axis.
Then $D_{n}$ is generated by a rotation $r$ through angle $\theta=\frac{2\pi}{n}$ around the $z$-axis, and a rotation $h$ through angle $\pi$ around the $x$-axis.
Specifically,
$$r=\left[\begin{matrix}
    \cos \theta & -\sin \theta &0\\
    \sin \theta & \cos \theta & 0\\
    0 & 0 & 1
\end{matrix}\right],\quad
h=\left[\begin{matrix}
    1 & 0 & 0\\
    0 & -1 & 0\\
    0 & 0 & -1
\end{matrix}\right].
$$
These matrices satisfy $r^n=I,h^2=I,hrh=r^{-1}.$
Therefore $D_n=\braket{r,h}=\{r^k,hr^k\mid k\in[n]\}.$
Let $L$ be the space spanned by the $z$-axis in $\mathbb{R}^3,$ and $P$ be the $x$-$y$ plane.
Notice that every element in $D_n$ is block diagonal relative to $P\oplus L.$
So $\mathrm{Span}_{\mathbb{R}}\{R_b\mid b\in B(f)\}\subseteq\mathrm{End}(P)\oplus \mathrm{End}(L)$.
By inspection, we can see $D_n$ generates every element in $\mathrm{End}(P)\oplus \mathrm{End}(L)$.
Thus, $\mathrm{Span}_{\mathbb{R}}\{R_b\mid b\in B(f)\}=\mathrm{End}(P)\oplus \mathrm{End}(L)$.
By \cref{lm: Cb and Rb}, $W=\mathrm{Span}_{\mathbb{R}}\{C(b)\mid b\in B(f)\}=\mathrm{End}(P)\oplus \mathrm{End}(L).$
Next we compute $W^\perp.$
Suppose $X\in W^\perp$, then $\mathrm{Tr}(C^{\tt T}X)=0$ for every $C\in \mathrm{End}(P)\oplus \mathrm{End}(L).$
$C$ has the form $\left[\begin{smallmatrix}
    A & 0\\
    0 & \alpha
\end{smallmatrix}\right]$ where $A\in M_2(\mathbb{R})$ and $\alpha\in \mathbb{R}$.
We write $X=\left[\begin{smallmatrix}
    Y & u\\
    v & \beta
\end{smallmatrix}\right]$ where $Y\in M_{2}(\mathbb{R})$, $u\in \mathbb{R}^{2\times 1}$, $v\in \mathbb{R}^{1\times 2}$, $\beta\in \mathbb{R}$.
Then 
$$
\mathrm{Tr}\left(
\left[\begin{matrix}
    A^{\tt T} & 0\\
    0 & \alpha
\end{matrix}\right]
\left[\begin{matrix}
    Y & u\\
    v & \beta
\end{matrix}\right]
\right)
=\mathrm{Tr}(A^{\tt T}Y)+\alpha\beta=0
$$
for every $A\in M_{2}(\mathbb{R})$ and $\alpha\in \mathbb{R}.$
Thus, $Y=O$ and $\beta=0.$
We have 
$$
W^{\perp}=\{\left[\begin{matrix}
    O & u\\
    v & 0
\end{matrix}\right]\mid u\in \mathbb{R}^{2\times 1}, v\in \mathbb{R}^{1\times 2}\}.
$$
By Pauli slice \cref{lm: Pauli slice}, $T^{(p)},T^{(q)}$ and $T^{(r)}\in W^\perp$ for every $p,q,r\in\{1,2,3\}.$
By inspection, the only possibility is that $t_{pqr}=0$ for every $p,q,r\in\{1,2,3\}.$
Thus, we prove the following lemma:
\begin{lemma}\label{lm: arity 6 dihedral group 3rd orth}
    If $B(f)\cong D_{n}(n\ge 3)$, then $f$ satisfies {\sc 3rd-Orth}.
\end{lemma}


\subsection{Klein Group}
Let $B(f)=\{b_0,b_1,b_2,b_3\}\cong K_4$.
It is proved in \cite{MingjiProgram} that we only need to consider the following two cases:
\begin{enumerate}
    \item
    (Standard Case).
    $$
    M(b_0)=\left[\begin{matrix}
    1 & 0\\
    0 & 1
    \end{matrix}\right],
    M(b_1)=\left[\begin{matrix}
    0 & \mathfrak{i}\\
    \mathfrak{i} & 0
    \end{matrix}\right],
    M(b_2)=\left[\begin{matrix}
    \mathfrak{i} & 0\\
    0 & -\mathfrak{i}
    \end{matrix}\right],
    M(b_3)=\left[\begin{matrix}
    0 & 1\\
    -1 & 0
    \end{matrix}\right].$$
        
    \item
    (Non-standard Case).
    $$M(b_0)=\left[\begin{matrix}
    1 & 0\\
    0 & 1
    \end{matrix}\right],
    M(b_1)=\left[\begin{matrix}
    \mathfrak{i} & 0\\
    0 & -\mathfrak{i}
    \end{matrix}\right],
    M(b_2)=\frac{1}{\sqrt{2}}\left[\begin{matrix}
    0 & 1+\mathfrak{i}\\
    -1+\mathfrak{i} & 0
    \end{matrix}\right],
    M(b_3)=\frac{1}{\sqrt{2}}\left[\begin{matrix}
    0 & -1+\mathfrak{i}\\
    1+\mathfrak{i} & 0
    \end{matrix}\right].$$
\end{enumerate}

\subsubsection{Standard Case}
By \cref{para: SU2 and SO3}, we can compute the corresponding rotation matrices of $b_0,b_1,b_2,b_3$ in $\mathbf{SO}(3)$:
$$
R_{b_0}=
\left[\begin{matrix}
    1 & 0 & 0\\
    0 & 1 & 0\\
    0 & 0 & 1\\
\end{matrix}\right],
R_{b_1}=
\left[\begin{matrix}
    1 & 0 & 0\\
    0 & -1 & 0\\
    0 & 0 & -1\\
\end{matrix}\right],
R_{b_2}=
\left[\begin{matrix}
    -1 & 0 & 0\\
    0 & -1 & 0\\
    0 & 0 & 1\\
\end{matrix}\right],
R_{b_3}=
\left[\begin{matrix}
    -1 & 0 & 0\\
    0 & 1 & 0\\
    0 & 0 & -1\\
\end{matrix}\right].
$$
So $\mathrm{Span}_{\mathbb R}\{R_{b}\mid b\in B(f)\}=\{\mathrm{diag}(x,y,z)\mid x,y,z\in \mathbb R\}.$
By \cref{lm: Cb and Rb}, $W:=\mathrm{Span}_{\mathbb R}\{C(b)\mid b\in B(f)\}=\{\mathrm{diag}(x,y,z)\mid x,y,z\in \mathbb R\}.$
So $W^\perp=\{A\in M_3(\mathbb{R})\mid A_{11}=A_{22}=A_{33}=0\}.$
By Pauli slice \cref{lm: Pauli slice}, $t_{pqr}\neq 0\implies p,q,r\text{ are mutually distinct}.$

In the standard case, by the realnumberizing method (see \cite{MingjiProgram} v3, Section 5), $f$ is a real signature up to a normalization.
Since normalization doesn't change the complexity, we assume $f$ is real.
Then $\rho_{ijk}=M(f)M(f)^{\dagger}=M(f)M(f)^{\tt T}=\rho_{ijk}^{\tt T}.$
It follows that $t_{pqr}=\mathrm{Tr}(\rho_{ijk}(\sigma_p\otimes \sigma_q\otimes \sigma_r))=\mathrm{Tr}((\sigma_p\otimes \sigma_q\otimes \sigma_r)^{\tt T}\rho_{ijk}^{\tt T})=\mathrm{Tr}(\rho_{ijk}(\sigma_p\otimes \sigma_q\otimes \sigma_r)^{\tt T})$.
Assume that $p,q,r$ are mutually distinct. 
Then $(p,q,r)$ contains exactly one $2$.
Notice that $\sigma_1^{\tt T}=\sigma_1,\sigma_2^{\tt T}=-\sigma_2,\sigma_3^{\tt T}=\sigma_3$, we have  $(\sigma_p\otimes \sigma_q\otimes \sigma_r^{\tt T})=-\sigma_p\otimes \sigma_q\otimes \sigma_r.$
So $t_{pqr}=\mathrm{Tr}(\rho_{ijk}(\sigma_p\otimes \sigma_q\otimes \sigma_r)^{\tt T})=-\mathrm{Tr}(\rho_{ijk}(\sigma_p\otimes \sigma_q\otimes \sigma_r))=-t_{pqr}$.
Thus, $t_{qpr}=0$ if $p,q,r$ are mutually distinct.
We have proved:
\begin{lemma}\label{lm: K4 standard case}
    In the standard case of $B(f)\cong K_4$, $f$ satisfies {\sc 3rd-Orth}.
\end{lemma}


\subsubsection{Non-standard Case}
In this case the Pauli slice constraints alone do not force all the
weight-three Pauli coefficients to vanish.  We instead use the full
factorization assertion in \cref{prop: binary closure}.

\begin{lemma}[Common matching]\label{lm: common matching}
    Up to a permutation of variables,
    \begin{equation}\label{eq: common matching}           
        f=\sum_{s=0}^3w_sb_s(x_1,x_2)b_{p(s)}(x_3,x_4)b_{q(s)}(x_5,x_6),
    \end{equation}
    where $w_s\neq 0\,(\forall s, 0\le s\le 3)$ and $p,q$ are both permutations of $\{0,1,2,3\}.$
\end{lemma}
\begin{proof}
    Suppose $f$ is normalized with norm 1.
    Define $u_i=\frac{1}{\sqrt{2}}({b_i}(00),b_i(01),b_i(10),b_i(11))^{\tt T}$ be the normalized vector.
    Notice that 
    $$\braket{u_i,u_j}=\frac{1}{2}\mathrm{Tr}(M(b_i)^\dagger M(b_j))=\delta_{ij},\quad  \forall i,j=0,1,2,3.$$ 
    So $\{u_0,u_1,u_2,u_3\}$ is an orthonormal basis for the vector space $ (\mathbb C^{4}).$
    Let $F=M_{x_1x_2,x_3x_4x_5x_6}(f)\in \mathbb{C}^{4\times 16}$, and $F_y$ denote the column in $F$ indexed by $y=(x_3,x_4,x_5,x_6).$ Then $F_y$ can be expanded as $F_y=\sum_{s=0}^3 c_s(y)u_s$, where $c_s(y)\in \mathbb{C}.$  
    Define $c_s=(c_s(0000),\ldots,c_s(1111))^{\tt T}\in \mathbb{C}^{16}.$
    Then $F=\sum_{s=0}^3 u_s c_s^{\tt T}$.
    The coefficient is $c_s(y)=u_s^\dagger F_y=\frac{1}{\sqrt{2}}\partial_{12}^{\overline{b_s}}f(y).$
    Since $f$ satisfies {\sc 2nd-Orth}, we have $$\frac{1}{4}\sum_{s=0}^3u_su_s^\dagger=\frac{1}{4}I_4=\rho_{12}=FF^\dagger=\sum_{s,t=0}^3u_sc_s^{\tt T}\overline{c_t}u_t^\dagger=\sum_{s,t=0}^3\braket{c_t,c_s}u_su_t^\dagger.$$
    Then $\braket{c_t,c_s}=\frac{1}{4}\delta_{st}$.
    Let $g_s=2c_s$, then $\braket{g_s,g_t}=\delta_{st}$, and $F=\frac{1}{2}\sum_{s=0}^3 u_sg_s^{\tt T}$.
    We can view $g_s$ as a signature on $y=(x_3,x_4,x_5,x_6).$
    Then $g_s=\sqrt{2}\partial_{12}^{\overline{b_s}}f$ and
    \begin{equation*}\label{eq: Schmidt decomposition of f}
        f(x_1,x_2,x_3,x_4,x_5,x_6)=\frac{1}{2}\sum_{s=0}^3 u_s(x_1,x_2)\otimes g_s(x_3,x_4,x_5,x_6).
    \end{equation*}
    Tracing out $x_1,x_2$, we have 
    \begin{equation}\label{eq: tracing out x1 x2}
        \rho_{3456}=\frac{1}{4}\sum_{s=0}^3g_sg_s^{\dagger}.
    \end{equation}
    Here \(g_sg_s^\dagger\) simply means the matrix whose \((x,y)\)-entry is \(g_s(x)\overline{g_s(y)}\).
    Since we can realize $g_s=\sqrt{2}\partial_{12}^{\overline{b_s}}f$ for $s=0,1,2,3$,
    by assumption \eqref{eq: arity 6 assumption}, $g_s$ factors into two binaries $b_{\ell_s}$ and $b_{k_s}.$
    The two binaries match up the four variables $\{x_3,x_4,x_5,x_6\}$ into two pairs.
    Consider the pair $\{3,4\}$.
    Let $k_{34}$ be the number of the four signatures \(g_0,g_1,g_2,g_3\) whose matching contains the edge \(34\).
    Tracing out $x_5,x_6$ in \eqref{eq: tracing out x1 x2},
    %We now consider the mating signature $\mathfrak{m}_{34}f$ (or the marginal of $\ket{f}$ on $x_3,x_4$).
    \begin{itemize}
        \item If $g_s=u_{\ell_s}(x_3,x_4)\otimes u_{k_s}(x_5,x_6)$, then the $s$-th term in \eqref{eq: tracing out x1 x2} contributes $\frac{1}{4}{u_{\ell_s}}{u_{\ell_s}}^\dagger.$
        
        \item If $g_s=u_{\ell_s}(x_3,x_5)\otimes u_{k_s}(x_4,x_6)$, then the $s$-th term in \eqref{eq: tracing out x1 x2} contributes $\frac{1}{4}(\frac{1}{2}I_2\otimes \frac{1}{2}I_2)$.
    \end{itemize}
    By {\sc 2nd-Orth}, $$\rho_{34}=\frac{1}{4}\left(\sum_{s: 34 \text{ is matched in } g_s}u_{\ell_s}u_{\ell_s}^\dagger+(4-k_{34})\frac{1}{4}I_4 \right)=\frac{1}{4}I_4.$$
    Multiplying by \(4\) and rearranging gives
    \begin{equation*}
        \sum_{s:\,34\text{ is matched in }g_s}u_{\ell_s}u_{\ell_s}^\dagger = \frac{k_{34}}4I_4.
    \end{equation*}
    Let \(m_j\) be the number of times the label \(u_j\) occurs on the edge \(34\).
    So
    $$\sum_{j=0}^3 m_j u_{j}u_j^\dagger=\frac{k_{34}}{4}I_4=\frac{k_{34}}{4} \sum_{j=0}^3 u_ju_j^\dagger.$$
    Multiply $u_i$ on both sides, we have $m_iu_i=\frac{k_{34}}{4}u_i,\,\forall i=0,1,2,3.$
    So $m_0=m_1=m_2=m_3=\frac{k_{34}}{4}.$
    Since $m_i\in \mathbb{N},$ we have $k_{34}=0$ or $4.$
    Moreover, if $k_{34}=4$, then $m_0=m_1=m_2=m_3=1.$
    Similarly, we can define $k_e$ be number of the four signatures $g_0,g_1,g_2,g_3$ whose matching contains the edge $e=ij$, and the same argument shows $k_e=0$ or $4.$

    Now we count the edge occurrences and show that $k_{34}=k_{56}=4$ after possibly renaming the variables.
    The four signatures \(g_0,g_1,g_2,g_3\) each have a perfect matching consisting of two edges. 
    Therefore the total number of edge occurrences is
    $4\cdot2=8.$
    For every residual edge \(e\), since $k_e=0$ or $4$, its number of occurrences is either \(0\) or \(4\). The only way these counts can sum to \(8\) is that exactly two edges occur four times each.
    Therefore all four contractions use the same matching.
    i.e., after renaming variables possibly, we may assume $$g_s(x_3,x_4,x_5,x_6)=u_{p(s)}(x_3,x_4)\otimes u_{q(s)}(x_5,x_6),\quad \text{for all } s=0,1,2,3.$$
    On either edge of that common matching, the factors \(u_0,u_1,u_2,u_3\) occur exactly once.
    That is, $p(s),q(s)$ are both permutations of $\{0,1,2,3\}.$
    The lemma is proven.
\end{proof}

\begin{lemma}
Under assumption \eqref{eq: arity 6 assumption}, the non-standard case of
$B(f)\cong K_4$ cannot occur.
\end{lemma}

\begin{proof}
    We note that $$M(b_r)M(b_s)=M(b_{r\oplus s}),$$
    where $\oplus$ is the commutative addition of the indices of the Klein group $B(f)$ defined by
    $$
    s\oplus s=0, \quad 1\oplus2=3, \quad 1\oplus 3=2,\quad 2\oplus 3=1.
    $$
    Also note that $B(f)$ has the transpose closure property:
    $$
    M(b_0)^{\tt T}=M(b_0), \quad
    M(b_1)^{\tt T}=M(b_1), \quad
    M(b_2)^{\tt T}=M(b_3), \quad 
    M(b_3)^{\tt T}=M(b_2).
    $$
    Let \(\tau(s)\) denote the label obtained by transpose. 
    Thus
    \[
    \tau(0)=0,\quad \tau(1)=1,\quad
    \tau(2)=3,\quad \tau(3)=2.
    \]
    Now we start from the common matching expression \eqref{eq: common matching} in \cref{lm: common matching}:
    \begin{equation*}    
        f=\sum_{s=0}^3w_sb_s(x_1,x_2)b_{p(s)}(x_3,x_4)b_{q(s)}(x_5,x_6),
    \end{equation*}
    Merge $x_2$ and $x_3$ using $=_2$, we get 
    \begin{equation}\label{eq: first expansion of g}
        \begin{aligned}
            \partial_{23}f(x_1,x_4,x_5,x_6)
            =&\sum_{s=0}^3\sum_{a=0}^1 w_sb_s(x_1,a)b_{p(s)}(a,x_4)b_{q(s)}(x_5,x_6)\\
            =&\sum_{s=0}^3w_sb_{s\oplus p(s)}(x_1,x_4)b_{q(s)}(x_5,x_6).
        \end{aligned}
    \end{equation}
    
    Next we show that $s\mapsto s\oplus p(s)$ is also a permutation of $\{0,1,2,3\}.$
    Since $\{b_0,b_1,b_2,b_3\}$ is a basis for the set of binary signatures, $\{b_r\otimes b_t\mid 0\le r,t\le 3\}$ is a basis for the set of 4-ary signatures.
    We can expand $\partial_{23}f$ under this basis:
    \begin{equation}\label{eq: second expansion of g}
        \partial_{23}f(x_1,x_4,x_5,x_6)=\sum_{r,t=0}^3 A_{rt}b_r(x_1,x_4)b_t(x_5,x_6).
    \end{equation}
    Compare \eqref{eq: first expansion of g} and \eqref{eq: second expansion of g}, we have $A_{s\oplus p(s),q(s)}=w_s.$
    Since $q(s)$ is a permutation, every column of $A$ receices a non-zero entry. 
    So $\mathrm{rank}(A)=\#\{s\oplus p(s)\mid s=0,1,2,3\}.$

    On the other hand, $\partial_{23}f$ factors into two binaries by assumption \eqref{eq: arity 6 assumption}.
    There are three cases:
    \begin{enumerate}
        \item $\partial_{23}f(x_1,x_4,x_5,x_6)=\lambda b_{i}(x_1,x_4)b_j(x_5,x_6)$.
        Then the matrix $A$ only has one non-zero entry $A_{ij}=\lambda.$
        So $\mathrm{rank}(A)=1.$
        \item $\partial_{23}f(x_1,x_4,x_5,x_6)=\lambda b_{i}(x_1,x_5)b_j(x_4,x_6)$.
        Then $M_{x_1x_4,x_5x_6}(\partial_{23}f)=\lambda M_{x_1,x_5}(b_i)\otimes M_{x_4,x_6}(b_j)$ and thus $\mathrm{rank}(M_{x_1x_4,x_5x_6}(\partial_{23}f))=2\times 2=4.$
        Basis change preserves the rank.
        So $\mathrm{rank}(A)=4.$
        \item $\partial_{23}f(x_1,x_4,x_5,x_6)=\lambda b_{i}(x_1,x_6)b_j(x_4,x_5)$.
        Similar to Case 2, $\mathrm{rank}(A)=4.$
    \end{enumerate}
    So far we proved $\mathrm{rank}(A)=1$ or $4$.
    If $\mathrm{rank}(A)=1,$ then $s\oplus p(s)=r$ for $s=0,1,2,3.$
    By \eqref{eq: first expansion of g}, $$\partial_{23}f=\sum_{s=0}^3w_sb_r(x_1,x_4)b_{q(s)}(x_5,x_6)=b_r(x_1,x_4)\otimes \sum_{s=0}^3w_sb_{q(s)}(x_5,x_6).$$
    But $q(s)$ is a permutation, \(b_0,b_1,b_2,b_3\) are linearly independent, and all the coefficients $w_s\neq 0$.
    Therefore $\sum_{s=0}^3 w_s b_{q(s)}$ cannot be proportional to a single member of \(B(f)\).
    This contradicts \cref{prop: binary closure}.
    Thus, $\mathrm{rank}(A)=\#\{s\oplus p(s)\mid s=0,1,2,3\}=4.$
    We proved that $s\mapsto s\oplus p(s)$ is a permutation of $\{0,1,2,3\}.$

    We now consider merging $x_1$ and $x_3$ using $=_2.$
    \begin{equation}\label{eq: merge 13}
        \begin{aligned}
            \partial_{13}f(x_2,x_4,x_5,x_6)
            =&\sum_{s=0}^3\sum_{a=0}^1 w_sb_s(a,x_2)b_{p(s)}(a,x_4)b_{q(s)}(x_5,x_6)\\
            =&\sum_{s=0}^3w_sb_{\tau(s)\oplus p(s)}(x_2,x_4)b_{q(s)}(x_5,x_6).
        \end{aligned}
    \end{equation}
    By the same rank argument above, $s\mapsto \tau(s)\oplus p(s)$ is a permutation of $\{0,1,2,3\}.$
    We are ready to arrive at a contradiction by the following claim:
    \begin{claim}\label{claim: impossible permutations}
        It is impossible that $p(s),s\oplus p(s),\tau(s)\oplus p(s)$ are all permutations of $\{0,1,2,3\}.$
    \end{claim}
    \begin{claimproof}{\cref{claim: impossible permutations}}
        Let
        $
        d=p(0), r(s)=p(s)\oplus d.
        $
        Since $p$ is a permutation, $r$ is also a
        permutation, and $r(0)=0$.
        By assumption, $s\mapsto s\oplus p(s)$ is a permutation. 
        If $r(s)=s$ for some $s\neq0$, then
        $
        s\oplus p(s)=s\oplus r(s)\oplus d=d=0\oplus p(0),
        $
        a contradiction. 
        Hence $r$ has no fixed point on $\{1,2,3\}$, so
        \[
        r|_{\{1,2,3\}}\in\{(1\,2\,3),(1\,3\,2)\}.
        \]
        In the first case, $r(2)=3=\tau(2)$; in the second case,
        $r(3)=2=\tau(3)$. Thus, in either case, there exists
        $s_*\in\{2,3\}$ such that $r(s_*)=\tau(s_*)$. Consequently,
        \[
        \tau(s_*)\oplus p(s_*)
         =\tau(s_*)\oplus r(s_*)\oplus d
         =d
         =\tau(0)\oplus p(0).
        \]
        Therefore $s\mapsto\tau(s)\oplus p(s)$ is not injective, hence is not
        a permutation. It is thus impossible for
        $
        p(s), s\oplus p(s), \tau(s)\oplus p(s)
        $
        to all be permutations of $\{0,1,2,3\}$.
    \end{claimproof}
    The lemma is proved.
\end{proof}


\subsection{Trivial Group}
In this subsection, we assume $B(f)\cong C_1$ is the trivial group, i.e., the only binary signature generated from $f$ is $=_2$.
For a pair of distinct indices $\{i,j\}$, we can view it as an edge in $K_6$.
We write $\partial_{ij}f=\partial_ef$.
\cref{prop: binary closure} gives
\begin{equation}
    \partial_e f
    =
    \lambda_e \bigotimes_{d\in M_e} (=_2)_d ,
    \qquad \lambda_e\neq 0,
    \label{eq:C1-contraction}
\end{equation}
where \(M_e\) is a perfect matching of the four vertices outside \(e\);
in particular, \(M_e\) consists of two disjoint edges.
We first show that all \(\lambda_e\) have the same modulus.  
By {\sc 2nd-Orth}, there is a common \(\mu>0\) such that
$
    |f_e^{ab}|^2=\mu,
    \text{for every \(e\) and \(a,b\in\{0,1\}\),}
$
and the four slices \(f_e^{ab}\) are pairwise orthogonal.  Hence
$
    |\partial_e f|^2
    =
    |f_e^{00}+f_e^{11}|^2
    =
    2\mu.
$
Since \(|(=_2)|^2=2\), the right-hand side of
\eqref{eq:C1-contraction} has squared norm \(4|\lambda_e|^2\).
Consequently,
$
    |\lambda_e|^2=\frac{\mu}{2},
$
which is independent of \(e\).

Now let \(e,d\) be disjoint edges.  The contractions commute:
$
    \partial_d\partial_e f
    =
    \partial_e\partial_d f.
$
If \(d\in M_e\), then the second contraction closes an equality loop
and contributes a factor \(2\).  If \(d\notin M_e\), then \(d\) crosses
the two matching edges and contributes no additional scalar.  In both
cases the remaining tensor is equality on the two uncontracted
vertices.  Therefore
\begin{equation}
    \lambda_e\,2^{[d\in M_e]}
    =
    \lambda_d\,2^{[e\in M_d]},
    \label{eq:C1-double-contraction}
\end{equation}
where \([P]\in\{0,1\}\) is the indicator of the statement \(P\).
Taking absolute values in \eqref{eq:C1-double-contraction} and using
\(|\lambda_e|=|\lambda_d|\) gives the reciprocity relation
\begin{equation}
    d\in M_e
    \quad\Longleftrightarrow\quad
    e\in M_d.
    \label{eq:C1-reciprocity}
\end{equation}
The powers of \(2\) in \eqref{eq:C1-double-contraction} then agree, so
\(\lambda_e=\lambda_d\) whenever \(e\cap d=\varnothing\).  
Hence there is
one common nonzero scalar
$
    \lambda_e=\lambda\neq 0
    \text{ for every edge \(e\).}
$
Define
$
    F_e:=\{e\}\cup M_e.
$
Then \(F_e\) is a perfect matching of \(K_6\).  
If \(d\in M_e\), write
\(M_e=\{d,c\}\).  By \eqref{eq:C1-reciprocity}, \(e\in M_d\), and the
remaining edge is necessarily \(c\).  Thus
$
    M_d=\{e,c\},
    \text{ and therefore }
    F_d=F_e.
$
It follows that the distinct sets \(F_e\) partition the \(15\) edges of
\(K_6\) into five perfect matchings.  
Up to relabelling, they are
\begin{equation}
\begin{aligned}
    &12\mid 34\mid 56, \qquad
     13\mid 25\mid 46, \qquad
     14\mid 26\mid 35,\\
    &15\mid 24\mid 36, \qquad
     16\mid 23\mid 45.
\end{aligned}
\label{eq:C1-factorization}
\end{equation}
We now derive a contradiction from four of these factors.  From
\[
    M_{12}=\{34,56\},\qquad
    M_{13}=\{25,46\},\qquad
    M_{25}=\{13,46\},\qquad
    M_{14}=\{26,35\},
\]
equation \eqref{eq:C1-contraction}, evaluated at suitable residual
assignments, gives
\begin{align}
    f_{001010}+f_{111010}&=0,
    &
    f_{000110}+f_{110110}&=0,
    \label{eq:C1-eight-a}\\
    f_{010010}+f_{111010}&=\lambda,
    &
    f_{000110}+f_{101110}&=0,
    \label{eq:C1-eight-b}\\
    f_{000000}+f_{010010}&=\lambda,
    &
    f_{100100}+f_{110110}&=0,
    \label{eq:C1-eight-c}\\
    f_{000000}+f_{100100}&=\lambda,
    &
    f_{001010}+f_{101110}&=\lambda.
    \label{eq:C1-eight-d}
\end{align}
The first seven equations imply
\[
\begin{aligned}
    f_{001010}&=-f_{111010},&
    f_{101110}&=f_{110110},\\
    f_{000000}&=f_{111010},&
    f_{100100}&=-f_{110110}.
\end{aligned}
\]
Using the first equation in \eqref{eq:C1-eight-d}, we therefore obtain
\[
\begin{aligned}
    f_{001010}+f_{101110}
    &=-f_{111010}+f_{110110}\\
    &=-\bigl(f_{000000}+f_{100100}\bigr)\\
    &=-\lambda.
\end{aligned}
\]
However, the second equation in \eqref{eq:C1-eight-d} says that the same
quantity is \(\lambda\).  Hence \(\lambda=-\lambda\), so
\(\lambda=0\), contradicting \(\lambda\neq 0\).
Therefore, we proved:
\begin{lemma}
    Under assumption \eqref{eq: arity 6 assumption}, the case $B(f)\cong C_1$ cannot occur.
\end{lemma}

\subsection{Order Two Group}
In this subsection, we assume that $B(f)\cong C_2$, the group of order $2$. 
We will show this case cannot occur under assumption \eqref{eq: arity 6 assumption}.
Suppose $B(f)=\{e,b\}$, where $M(e)=I_2$ and $b^2=e.$
\begin{lemma}\label{lm: C2 standard form}
    Up to a normalization and a real orthogonal holographic transformation, there are two cases: 1. (Symmetric Case): $M(b)=\left[\begin{matrix}
        \mathfrak{i} & 0\\
        0 & -\mathfrak{i}
    \end{matrix}\right]$
    or 2. (Asymmetric Case):
    $M(b)=\left[\begin{matrix}
        0 & 1\\
        -1 & 0
    \end{matrix}\right].$
\end{lemma}
\begin{proof}
    Notice that $B(f)$ is closed under transpose. 
    So $M(b)^{\tt T}=\pm M(b).$
    \begin{enumerate}
        \item $M(b)^{\tt T}=M(b)$.
        Assume $M(b)=\left[\begin{matrix}
            x & y\\
            y & z
        \end{matrix}\right]\in \mathbf{SU}(2).$
        Since $b$ has order two,
        we have $M(b)^2=\left[\begin{matrix}
            x^2+y^2 & xy+yz\\
            xy+yz & y^2+z^2
        \end{matrix}\right]=\pm I.$
        So $y(x+z)=0$ and $x^2=z^2.$
        \begin{itemize}
            \item 
            If $y=0$, $M(b)=\left[\begin{matrix}
            x & 0\\
            0 & z
            \end{matrix}\right].$
            By $\det(M(b))=xz=1$ and $x^2=z^2$, we have $x^4=1$.
            If $x=\pm 1$, then $z=\pm 1$ by $xz=1$, contradicting $b$ has order 2.
            So $x=\pm \mathfrak{i}.$
            It follows that $z=\mp \mathfrak{i}.$
            Up to a normalization, $M(b)=\left[\begin{matrix}
                \mathfrak{i} & 0\\
                0 & -\mathfrak{i}
            \end{matrix}\right].$

            \item If $x+z=0$, $M(b)=\left[\begin{matrix}
                x & y\\
                y & -x
            \end{matrix}\right].$
            By $M(b)\in \mathbf{SU}(2)$, 
            we have $x^2+y^2=-1$ and 
            $M(b)M(b)^\dagger=\left[\begin{matrix}
                |x|^2+|y|^2 & x\overline{y}-\overline{x}y\\
                \overline{x}y-x\overline{y} & |x|^2+|y|^2
            \end{matrix}\right]=\pm I_2$. 
            So $x\overline{y}=\overline{x}y$.
            If $xy\neq 0$, then $\arg(x)=\arg(y)$.
            After a normalization, we may assume $x,y\in \mathbb{R}$.
            If $xy=0$, we may still normalize and assume $x,y\in \mathbb{R}.$
            Since $x^2+y^2=-1$, we may assume $x=\mathfrak{i}\cos(2\theta)$ and $y=\mathfrak{i}\sin(2\theta).$
            Let $T=\left[\begin{matrix}
                \cos \theta & -\sin \theta\\
                \sin \theta & \cos \theta
            \end{matrix}\right]\in \mathbf{SO}(2).$
            A direct calculation shows that $T^{\tt T}M(b)T=\left[\begin{matrix}
                \mathfrak{i} & 0\\
                0 & -\mathfrak{i}
            \end{matrix}\right].$
            So the lemma is true by a $T$-transformation.
        \end{itemize}
        
        \item 
        $M(b)^{\tt T}=-M(b).$
        Assume that 
        $M(b)=\left[\begin{matrix}
                0 & x\\
                -x & 0
            \end{matrix}\right].$
        Since $M(b)\in \mathbf{SU}(2)$, we have $x^2=1$, $x=\pm 1$.
        So up to a normalization, $M(b)=\left[\begin{matrix}
                0 & 1\\
                -1 & 0
            \end{matrix}\right].$
    \end{enumerate}
\end{proof}

For simplicity, in the standard case, we no longer require $M(b)\in \mathbf{SU}(2)$ 
and assume that $M(b)=\left[\begin{matrix}
    1 & 0\\
    0 & -1
\end{matrix}\right].$
We first record two lemmas that will be used in both cases:
\begin{lemma}\label{lem:c2-two-projection}
    Suppose that $g=\lambda\phi(x_\ell,x_m)\psi(x_u,x_v)$ for some $\lambda\neq 0$, $\phi,\psi\in\{e,b\}$ and $\{1,2\}\subseteq\{\ell,m,p,q\}.$
    If $\{1,2\}=\{\ell,m\}$ or $\{1,2\}=\{p,q\}$, then exactly one of
    $\partial_{12}^{e}g$ and $\partial_{12}^{b}g$ is nonzero.  
    Otherwise, both $\partial_{12}^{e}g$ and $\partial_{12}^{b}g$ are nonzero and are not proportional.
\end{lemma}

\begin{proof}
    Define $\phi^p=e$ if $p \equiv 0 \text{ (mod }{2})$, and $\phi^p=b$ if $p \equiv 1 \text{ (mod }{2})$.
    First suppose that $g=\lambda\phi^p(x_1,x_2)\phi^q(x_u,x_v)$.
    Notice that
    \[
     \sum_{\alpha,\beta\in\{0,1\}}
     \phi^p(\alpha,\beta)\phi^q(\alpha,\beta)
     =
     \begin{cases}
     2,&p \equiv q \text{ (mod }{2})\\
     0,&p \not\equiv q \text{ (mod }{2}).
     \end{cases}
    \]
    Consequently, merging $x_1$ and $x_2$ by $e$ or $b$ gives exactly one nonzero result.

    Now suppose that $g=\lambda\phi^p(x_1,x_u)\phi^q(x_2,x_v)$.
    In the symmetric case, merging $x_1$ and $x_2$ by $\phi^t$ produces $\partial_{12}^{\phi^t} g=\lambda\phi^{p+q+t}(x_u,x_v)$.
    The two signatures $\partial_{12}^eg$ and $\partial_{12}^bg$ therefore produce the two distinct factors $e$ and $b$, up to nonzero scalars.
    In the asymmetric case, $M_{x_u,x_v}(\partial_{12}^{\phi^t} g)=\lambda M(\phi^p)^{\tt T}M(\phi^t)M(\phi^q).$
    For $t=0$ and $t=1$, these matrices are projectively $I_2$ and $M(b)$. 
    Hence they are again nonzero and not proportional.
\end{proof}

\begin{lemma}\label{lm: C2 contraction orth}
    For any pair of indices $\{i,j\},$ $|\partial_{ij}^{e}f|=|\partial_{ij}^bf|> 0$ and $\braket{\partial_{ij}^ef,\partial_{ij}^bf}=0.$
\end{lemma}

\begin{proof}
    In the symmetric case, $\partial_{ij}^ef=f_{ij}^{00}+f_{ij}^{11}$, and $\partial_{ij}^bf=f_{ij}^{00}-f_{ij}^{11}$.
    By {\sc 2nd-Orth}, $|f_{ij}^{00}|=|f_{ij}^{11}|>0$ and $\braket{f_{ij}^{00},f_{ij}^{11}}=0.$
    So
    \begin{equation*}
        \begin{aligned}
            |\partial_{ij}^{e}f|^2=\braket{f_{ij}^{00}+f_{ij}^{11},f_{ij}^{00}+f_{ij}^{11}}=|f_{ij}^{00}|^2+|f_{ij}^{11}|^2>0.
        \end{aligned}
    \end{equation*}
    Similarly, $|\partial_{ij}^{b}f|^2=|f_{ij}^{00}|^2+|f_{ij}^{11}|^2=|\partial_{ij}^{e}f|^2>0.$
    Also,
    \begin{equation*}
        \begin{aligned}
            \braket{\partial_{ij}^{e}f,\partial_{ij}^{b}f}=\braket{f_{ij}^{00}+f_{ij}^{11},f_{ij}^{00}-f_{ij}^{11}}=|f_{ij}^{00}|^2-|f_{ij}^{11}|^2-\braket{f_{ij}^{00},f_{ij}^{11}}+\braket{f_{ij}^{11},f_{ij}^{00}}
            =0.
        \end{aligned}
    \end{equation*}
    In the asymmetric case, $\partial_{ij}^ef=f_{ij}^{01}+f_{ij}^{10}$, and $\partial_{ij}^bf=f_{ij}^{01}-f_{ij}^{10}$.
    {\sc 2nd-Orth} gives $|f_{ij}^{01}|=|f_{ij}^{10}|>0$ and $\braket{f_{ij}^{01},f_{ij}^{10}}=0.$
    So $|\partial_{ij}^{e}f|^2=|\partial_{ij}^bf|^2=|f_{ij}^{01}|^2+|f_{ij}^{10}|^2> 0$ and $\braket{\partial_{ij}^ef,\partial_{ij}^bf}=0.$
\end{proof}

\subsubsection{Symmetric Case}

\todo{ZZT: I will replace the symmetric Case by a simpler argument.}
In this case, $B(f)=\{e,b\}$ where $M(e)=I_2$ and $M(b)=\left[\begin{matrix}
    1 & 0\\
    0 & -1
\end{matrix}\right]$.

\iffalse
\begin{lemma}\label{lm: C2 even support}
    For any $x\in\mathscr{S}(f)$, $\mathrm{wt}(x)$ is even.
\end{lemma}
\begin{proof}
    Suppose $a=(a_1,a_2,a_3,a_4,a_5,a_6)\in \mathscr{S}(f)$ but $\mathrm{wt}(a)$ is odd.
    After renaming the variables, we may assume $a_1=a_2,a_3=a_4,a_5\neq a_6.$
    Consider merging $x_1$ and $x_2$ using $b^s\in B(f)$, and merging $x_3$ and $x_4$ using $b^t\in B(f)$.
    The resulting signature $\partial_{12}^{b^s}\partial^{b^t}_{34}f$ is also in $B(f)$ by assumption \eqref{eq: arity 6 assumption}.
    So $\partial_{12}^{b^s}\partial^{b^t}_{34}f(a_5,a_6)=0.$
    Thus,
    \begin{equation}\label{eq: C2 even support}
        \sum_{i,j\in\{0,1\}}(-1)^{si+tj}f(i,i,j,j,a_5,a_6)=\partial_{12}^{b^s}\partial^{b^t}_{34}f(a_5,a_6)=0.
    \end{equation}
    Pick $(s,t)=(0,0),(0,1),(1,0)$ and $(1,1)$ respectively, \eqref{eq: C2 even support} gives that 
    \begin{equation*}
        \left[\begin{matrix}
            1 & 1 & 1 & 1\\
            1 & -1 & 1 & -1\\
            1 & 1 & -1 & -1\\
            1 & -1 & -1 & 1\\
        \end{matrix}\right]
        \left[\begin{matrix}
            f(0,0,0,0,a_5,a_6)\\
            f(0,0,1,1,a_5,a_6)\\
            f(1,1,0,0,a_5,a_6)\\
            f(1,1,1,1,a_5,a_6)\\
        \end{matrix}\right]
        =\left[\begin{matrix}
            0\\
            0\\
            0\\
            0\\
        \end{matrix}\right].
    \end{equation*}
    The coefficient matrix is 
    $\left[\begin{smallmatrix}
        1 & 1\\
        1 & -1
    \end{smallmatrix}\right]\otimes\left[\begin{smallmatrix}
        1 & 1\\
        1 & -1
    \end{smallmatrix}\right]$, which has rank 4.
    Thus $f(a)=0$, contradiction.
\end{proof}
\fi


\begin{lemma}\label{lm: C2 symmetric impossible}
    Under assumption \eqref{eq: arity 6 assumption}, the symmetric case of $B(f)\cong C_2$ cannot occur.
\end{lemma}
\begin{proof}
    
    By permuting the residual variables $(x_3,x_4,x_5,x_6)$ and applying local binary modification by $b$, we may normalize
    \begin{equation}
         \partial_{12}^{e}f=x\cdot e(x_3,x_4)e(x_5,x_6),
         \qquad x\in\mathbb C^\times.
         \label{eq:c2-symmetric-normalization}
    \end{equation}
    Here $x\neq 0$ by \cref{lm: C2 contraction orth}.
    With a little abuse of notation, we still use $f$ to denote the 6-ary signature after the binary modification.
    Notice that a binary modification by $b$ doesn't change assumption \eqref{eq: arity 6 assumption} and the conclusion in \cref{lem:c2-two-projection}.
    \begin{claim}\label{claim: C2 symmetry}
        Up to the symmetries preserving \eqref{eq:c2-symmetric-normalization}, there are three possibilities for $\partial_{12}^{b}f$:
        \[
         y\cdot e(x_3,x_4)b(x_5,x_6),\quad
         y\cdot b(x_3,x_4)b(x_5,x_6),\quad
         y\cdot e(x_3,x_5)b(x_4,x_6),
         \qquad y\in\mathbb C^\times.
        \]
    \end{claim}
    \begin{claimproof}{\cref{claim: C2 symmetry}}
        By \cref{lm: C2 contraction orth}, $\braket{\partial_{12}^{e}f,\partial_{12}^{b}f}=0.$
        If $\partial_{12}^bf=y\cdot \phi^p(x_3,x_4)\phi^q(x_5,x_6)$ where $y\neq 0$ by \cref{lm: C2 contraction orth}, orthogonality to \(\partial_{12}^{e}f=x\cdot e(x_3,x_4)e(x_5,x_6)\) requires at least one \(b\)-factor in $\phi^p$ and $\phi^q$. 
        Up to exchanging $(x_3,x_4)$ and $(x_5,x_6)$, we have $\partial_{12}^bf=y\cdot e(x_3,x_4)b(x_5,x_6)$ or $\partial_{12}^bf= y\cdot b(x_3,x_4)b(x_5,x_6).$
        Otherwise, $x_3$ and $x_4$ are in distinct binary factors of $\partial_{12}^bf$.
        Up to exchanging $x_5$ and $x_6$, we may assume
        $\partial_{12}^bf=y\cdot \phi^p(x_3,x_5)\phi^q(x_4,x_6), \,y\ne 0.$
        Then $0=\braket{\partial_{12}^ef,\partial_{12}^bf}=xy\cdot\mathrm{Tr}(I_2\cdot M(\phi^q)\cdot I_2\cdot M(\phi^p)^{\tt T})=(-1)^{p+q}.$
        So exactly one binary factor of $\partial_{12}^bf$ is $b$.
        Up to exchanging $(x_3,x_5)$ and $(x_4,x_6)$, $\partial_{12}^bf=y\cdot e(x_3,x_5)b(x_4,x_6).$
    \end{claimproof}
    
    
    Suppose first that
    $\partial_{12}^bf=y\cdot e(x_3,x_4)b(x_5,x_6)$, and put $g=\partial_{34}^{b}f$.  Commutativity of contractions gives
    \[
     \partial_{12}^{e}g
     =\partial_{34}^{b}\partial_{12}^{e}f=0,
     \qquad
     \partial_{12}^{b}g
     =\partial_{34}^{b}\partial_{12}^{b}f=0.
    \]
    Since $g\neq 0$ by \cref{lm: C2 contraction orth}, the above equations contradicts \cref{lem:c2-two-projection}.
    Next suppose that $\partial_{12}^{b}f=y\cdot b(x_3,x_4)b(x_5,x_6)$, and put $g=\partial_{35}^{e}f$.  Then
    \[
     \partial_{12}^{e}g
     =\partial_{35}^{e}\partial_{12}^{e}f=x\cdot e(x_4,x_6),
     \qquad
     \partial_{12}^{b}g
     =\partial_{35}^{e}\partial_{12}^{b}f=y\cdot e(x_4,x_6).
    \]
    Both projections are nonzero and proportional, again contradicting
    \cref{lem:c2-two-projection}.

    It remains to consider the case
    \begin{equation}
        \partial_{12}^{e}f=x\cdot e(x_3,x_4)e(x_5,x_6),
         \qquad
         \partial_{12}^{b}f=y\cdot e(x_3,x_5)b(x_4,x_6).
         \label{eq:c2-symmetric-crossed}
    \end{equation}
    Since $\partial_{12}^{e}f=f_{12}^{00}+f_{12}^{11}$ and $\partial_{12}^{b}f=f_{12}^{00}-f_{12}^{11},$
    we can recover $f_{12}^{00}$ and $f_{12}^{11}$ by 
    $f_{12}^{00}=\frac{1}{2}(\partial_{12}^{e}f+\partial_{12}^{b}f)$ and  $f_{12}^{11}=\frac{1}{2}(\partial_{12}^{e}f-\partial_{12}^{b}f).$
    In particular, we can list the truth table of $2f$ with \(x_1=x_2=x_3\) by \eqref{eq:c2-symmetric-crossed}:
    \begin{equation}\label{eq: C2 table 1=2=3}
        \begin{array}{c|cccccccc}
     \sigma&
     000000&000011&000101&000110&
     111001&111010&111100&111111\\ \hline
     2f^\sigma&
     x+y&x&-y&0&0&-y&x&x+y
    \end{array}.
    \end{equation}
    
    We next apply consider merging $x_1$ and $x_3$.
    Order the residual variables
    as $(x_2,x_4,x_5,x_6)$ and define
    \[
    \begin{aligned}
     R_{0}^{p,q}&=\phi^p(x_2,x_4)\phi^q({x_5,x_6}),\\
     R_{1}^{p,q}&=\phi_p(x_2,x_5)\phi^q(x_4,x_6),\\
     R_{2}^{p,q}&=\phi^p(x_2,x_6)\phi^q(x_4,x_5).
    \end{aligned}
    \]
    We have
    $\partial_{13}^{e}f=z\cdot R_{m}^{p,q}$ and
    $\partial_{13}^{b}f=w\cdot R_{m'}^{r,s}$ for some
    $z,w\in\mathbb C^\times$.
    Let $V_0(R)$ denote the values of $R(x_2,x_4,x_5,x_6)$ on
    \(
     (0000,0011,0101,0110),
    \)
    and let $V_1(R)$ denote the value of $R(x_2,x_4,x_5,x_6)$ on
    \(
     (1001,1010,1100,1111).
    \)
    Direct evaluation gives
    \[
    \begin{array}{c|c|c}
     m&V_0(R_{m}^{p,q})&V_1(R_{m'}^{r,s})\\ \hline
     0&(1,(-1)^q,0,0)&(0,0,(-1)^r,(-1)^{r+s})\\
     1&(1,0,(-1)^q,0)&(0,(-1)^r,0,(-1)^{r+s})\\
     2&(1,0,0,(-1)^q)&((-1)^r,0,0,(-1)^{r+s})
    \end{array}.
    \]
    Comparison with \eqref{eq: C2 table 1=2=3} gives
    \begin{align}
     zV_0(R_{m}^{p,q})+wV_0(R_{m'}^{r,s})
       &=(x+y,x,-y,0), \label{eq:c2-V0}\\
     zV_1(R_{m}^{p,q})-wV_1(R_{m'}^{r,s})
       &=(0,-y,x,x+y). \label{eq:c2-V1}
    \end{align}
    Because $x,y\ne0$, the second and third coordinates on the right side of \eqref{eq:c2-V0} are nonzero.  
    Hence $\{m,m'\}=\{0,1\}.$
    Solving the remaining signs in \eqref{eq:c2-V0}--\eqref{eq:c2-V1} gives exactly the following four possibilities:
    \[
    \begin{array}{c|c|c}
         \partial_{13}^{e}f&\partial_{13}^{b}f&
         \text{relations among the scalars}\\ \hline
         zR_{0}^{0,0}&wR_{1}^{0,1}&z=x,\ w=y\\
         zR_{0}^{1,1}&wR_{1}^{1,0}&x=w,\ y=-w,\ z=-w\\
         zR_{1}^{0,0}&wR_{0}^{0,1}&x=-w,\ y=w,\ z=-w\\
         zR_{1}^{1,1}&wR_{0}^{1,0}&x=w,\ y=z
    \end{array}.
    \]

    Now evaluate $\partial_{14}^{e}f$ on the residual words
    \(
     \mathcal W=(1010,0110,1001,0101),
    \)
    where the residual variables are ordered as $(x_2,x_3,x_5,x_6)$.  
    Using \eqref{eq:c2-symmetric-crossed} and the four rows above, respectively,
    we obtain
    \[
    \left.\partial_{14}^{e}f\right|_{\mathcal W}
     =
    \frac12(w,w,w,w),\text{ or }
    \frac12(-w,-w,-w,-w),\text{ or }
    \frac12(-w,w,w,-w),\text{ or }
    \frac12(-z,z,z,-z).
    \]
    As a sample of calculation, in the first row consider \((x_2,x_3,x_5,x_6)=1010\).
    Then $\partial_{14}^ef=f({010010})+f({110110}).$
    We have $$f({010010})=\frac{1}{2}(\partial_{13}^ef(1010)+\partial_{13}^b(1010))=\frac{1}{2}(zR_0^{0,0}(1010)+wR_1^{0,1}(1010))=\frac{w}{2},$$
    and 
    $$f({110110})=\frac{1}{2}(\partial_{12}^ef(0110)-\partial_{12}^b(0110))=0.$$
    So $\partial_{14}^ef(1010)=\frac{w}{2}.$
    The other entries are obtained in exactly the same way from $\partial_{12}^tf$ and $\partial_{13}^tf$, where $t\in\{e,b\}$.
    Thus $\partial_{14}^{e}f$ is nonzero on all four words of $\mathcal W$.
    On the other hand, $\partial_{14}^{e}f$ is a product of two diagonal binary signatures. 
    The intersection $\mathscr{S}(\partial_{14}^ef)\cap \mathcal W$, according two the matching of $x_2,x_3,x_5,x_6$, is
    \[
    \begin{array}{c|c}
     23\mid56&\varnothing\\
     25\mid36&\{1010,0101\}\\
     26\mid35&\{0110,1001\}
    \end{array}.
    \]
    Hence $\partial_{14}^ef$ is nonzero on at most two words of $\mathcal W$, contradiction.
    Therefore, the case \eqref{eq:c2-symmetric-crossed} cannot occur.
\end{proof}


\subsubsection{Asymmetric Case}
In the asymmetric case, we have $B(f)=\{e,b\}$, where $M(e)=I_2$ and $M(b)=\left[\begin{matrix}
    0 & 1\\
    -1 & 0
\end{matrix}\right].$
\begin{lemma}\label{lm: C2 asymmetric impossible}
    Under assumption \eqref{eq: arity 6 assumption}, the asymmetric case of $B(f)\cong C_2$ cannot occur.
\end{lemma}
\begin{proof}
    Similar to the proof of \cref{lm: C2 symmetric impossible}, we may permute the variables and apply binary modification and assume
    \[
     \partial_{12}^{e}f=x\cdot e(x_3,x_4)e(x_5,x_6),
     \qquad x\in\mathbb C^\times.
    \] 
    By \cref{lm: C2 contraction orth}, $\partial_{ij}^{e}f\neq 0,\,\partial_{ij}^bf\neq 0,$ and $\braket{\partial_{ij}^ef,\partial_{ij}^bf}=0$ for any pair of indices ${i,j}.$ 
    Again orthogonality of $\partial_{12}^{e}f$ and
    $\partial_{12}^{b}f$ leaves three possibilities for $\partial_{12}^bf:$
    \[
     y\cdot e(x_3,x_4)b(x_5,x_6),\quad
     y\cdot b(x_3,x_4)b(x_5,x_6),\quad
     y\cdot e(x_3,x_5)b(x_4,x_6),
    \]
    for some $y\in \mathbb{C}^\times.$
    Suppose first that $\partial_{12}^{b}f=y\cdot e(x_3,x_4)b(x_5,x_6)$, and put $g=\partial_{34}^{b}f$.  Commutativity gives
    \[
     \partial_{12}^{e}g
     =\partial_{34}^{b}\partial_{12}^{e}f=0,
     \qquad
     \partial_{12}^{b}g
     =\partial_{34}^{b}\partial_{12}^{b}f=0.
    \]
    This contradicts \cref{lem:c2-two-projection}.
    Next suppose that $\partial_{12}^bf=y\cdot b(x_3,x_4)b(x_5,x_6)$, and put $g=\partial_{35}^{e}f$.  Since $M(b)^{\tt T}I_2M(b)=I_2$, we have
    \[
     \partial_{12}^{e}g=\partial_{35}^e\partial_{12}^ef=x\cdot e(x_4,x_6),
     \qquad
     \partial_{12}^{b}g=\partial_{35}^e\partial_{12}^bf=y\cdot e(x_4,x_6).
    \]
    These are nonzero and proportional, again contradicting \cref{lem:c2-two-projection}.

    It remains to consider the case
    \begin{equation}
        \partial_{12}^{e}f=x\cdot e(x_3,x_4)e(x_5,x_6),
        \qquad     \partial_{12}^{b}f=y\cdot e(x_3,x_5)b(x_4,x_6).
    \end{equation}
    Put $g=\partial_{34}^{b}f$.  Commuting the two contractions gives
    \begin{equation}\label{eq: C2 asym}
        \partial_{12}^{e}g=\partial_{34}^b\partial_{12}^ef=0,
     \qquad
     \partial_{12}^{b}g=\partial_{34}^b\partial_{12}^bf=-y\cdot e(x_5,x_6)\ne0,
    \end{equation}
    where the second identity follows from $I_2^{\tt T}JJ=-I_2$.

    By \cref{lem:c2-two-projection}, $e(x_1,x_2)\mid g$ or $b(x_1,x_2)\mid g$.
    If $g=z\cdot e(x_1,x_2)t(x_5,x_6)$ for some $t\in \{e,b\}$, then $\partial_{12}^e g=2z\cdot t(x_5,x_6)\neq 0$, contradiction.
    Thus
    $g=z\cdot b(x_1,x_2)t(x_5,x_6)$ for some $t\in\{e,b\}$.  
    Merge $x_1$ and $x_2$ by $b$ gives
    \(
     \partial_{12}^{b}g=2z\cdot t(x_5,x_6).
    \)
    Comparison with \eqref{eq: C2 asym} gives $|y|=2|z|$ and $t=e$.
    On the other hand, for any pair of indices $\{i,j\},$
   
    $$|\partial_{ij}^{b}f|^2
    =
    \sum_{r,s,p,q}
    b(r,s)\overline{b(p,q)}
    \langle f_{ij}^{rs},f_{ij}^{pq}\rangle=\sum_{r,s}|b(r,s)|^2\braket{f_{ij}^{rs},f_{ij}^{rs}}=2\lambda$$
    by {\sc 2nd-Orth}, where $\lambda=|f_{ij}^{00}|^2=|f_{ij}^{01}|^2=|f_{ij}^{10}|^2=|f_{ij}^{11}|^2$.
    Consequently $|\partial_{12}^{b}f|=|\partial_{34}^{b}f|=|g|$.
    So $|y|=|z|$, contradicting $|y|=2|z|$.
\end{proof}

\subsection{Cyclic Group of Order at Least Three}
In this subsection, we assume that $B(f)\cong C_n(n\ge 3)$ which is the cyclic group of order $n$.
We are going to prove that under assumption \eqref{eq: arity 6 assumption}, $B(f)\cong C_n$ cannot occur.
In \cite{MingjiProgram} v1, it is proved that we only need to consider two possible cases:
\begin{enumerate}
    \item 
    (Standard Case).
    In the setting of $\Holant(\mathcal{F}),$ up to a real orthogonal holographic transformation, 
    the projective binary group of $f$ is $B(f)=\{b^0, b^1,\ldots,b^{n-1}\}$, where $b$ is the binary signature $\left[\begin{matrix}
        1 & 0\\
        0 & \zeta
    \end{matrix}\right]$ and  $\zeta=e^{\frac{2\pi \mathfrak{i}}{n}}.$
    
    \item 
    (Non-standard Case).
    In the setting of $\holant{\neq_2}{\hF}$, the projective binary group of $\widehat{f}$ is $\widehat{B}(\widehat{f})=\{b^0,b^1,\ldots,b^n\}$,
    where $b$ is the binary signature $\left[\begin{matrix}
        0 & 1\\
        \zeta & 0
    \end{matrix}\right]$ and $\zeta=e^{\frac{2\pi \mathfrak{i}}{n}}.$
    %\todo{$\widehat{B}(\widehat{f})$ has not been defined.}
\end{enumerate}

\subsubsection{Standard Case}
By \cref{prop: binary closure}, after merging $x_i$ and $x_j$ using some binary signature $b^k$ in $B(f)$, $\partial_{ij}^{b^k}f$ factors into two binary signatures in $B(f)$. 
This two binaries matches up the remaining variables by pairs.
The following lemma shows that for any $\{i,j\}$ there is a common matching which is independent of $k\in [n]$.
\begin{lemma}[Common matching]\label{lm: common matching Cn}
    Fix an arbitrary pair of indices $\{i,j\}$.
    For any $b^k\in B(f)$, $\partial_{ij}^{b^k}f(x_\ell,x_m,x_p,x_q)$ factors into binaries $b^{k_1}(x_\ell,x_n)$ and $b^{k_2}(x_p,x_q)$ for some $k_1,k_2\in [n]$ after permuting the variables possibly, where $\{i,j,\ell,m,p,q\}=\{1,2,3,4,5,6\}$.
\end{lemma}
\begin{proof}
    Let $f^{(k)}=\partial_{ij}^{b^k}f,\, g=f_{ij}^{00},\,h=f_{ij}^{11}$ and $y=(x_\ell,x_m,x_p,x_q).$
    We have 
    \begin{equation}\label{eq: cyclic group merging}
        f^{(k)}(y)=\partial_{ij}^{b^{k}}f(y)=f_{ij}^{00}(y)+\zeta^k f_{ij}^{11}(y)=g(y)+\zeta^k h(y), \quad \text{for every } y\in \{0,1\}^4.
    \end{equation}
    By assumption \eqref{eq: arity 6 assumption}, every merging signature $f^{(k)}$ factors into two binary signatures in $B(f).$
    %Without of loss of generality, we may assume $\{i,j\}=\{5,6\}.$
    In the standard case, the two binary signatures match up $x_\ell,x_m,x_p,x_q$ into two equal pairs.
    So the support of $f^{(k)}$ is
    \begin{equation*}
        \begin{aligned}
            &S_1=\{0000,0011,1100,1111\}, \text{ or }\\
            &S_2=\{0000,0101,1010,1111\}, \text{ or }\\
    &S_3=\{0000,0110,1001,1111\},
        \end{aligned}
    \end{equation*}
    according to the matching $\ell m\mid pq, \ell p\mid mq$ or $\ell q\mid mp.$
    \begin{claim}\label{claim: cyclic group common matching}
        If there exists two distinct $k,k'\in [n]$, such that $\mathscr{S}(f^{(k)})=\mathscr{S}(f^{(k')})=S$, then $\mathscr{S}(f^{(t)})=S$ for every $t\in[n]$. 
    \end{claim}
    \begin{claimproof}{\cref{claim: cyclic group common matching}}
    By \eqref{eq: cyclic group merging}, we have
    \[\left[
    \begin{matrix}
        f^{(k)}(y)\\
        f^{(k')}(y)
    \end{matrix}
    \right]
    =\left[\begin{matrix}
        1 & \zeta^k\\
        1 & \zeta^{k'}
    \end{matrix}
    \right]
    \left[\begin{matrix}
        g(y)\\
        h(y)
    \end{matrix}
    \right],\quad \forall y\in \{0,1\}^4.
    \]
    Thus, $g(y)=h(y)=0 \iff f^{(k)}(y)=f^{(k')}(y)=0.$
    %This implies $\mathscr{S}(f^{(k)})\cup \mathscr{S}(f^{(k')})=\mathscr{S}(g)\cup \mathscr{S}(h).$
    Since $\mathscr{S}(f^{(k)})=\mathscr{S}(f^{(k')})=S$, $g$ and $h$ vanishes outside $S$.
    By \eqref{eq: cyclic group merging}, $f^{(t)}$ vanishes outside $S$ for every $t\in [n]$.
    So $\mathscr{S}(f^{(t)})\subseteq S.$
    Since $\mathscr{S}(f^{(t)})=S_1,S_2$ or $S_3$, we have $\mathscr{S}(f)^{(t)}=S,$ for every $t\in [n].$
    \end{claimproof}
    
    If $n\ge 4$, then by pigeon hole principle, there exists distinct $k,k'\in [n]$ such that $\mathscr{S}(f^{(k)})=\mathscr{S}(f^{(k')})=S$.
    By \cref{claim: cyclic group common matching}, the lemma is proved.
    In the following, we may assume that $n=3$, and $\{\mathscr{S}(f^{(0)}),\mathscr{S}(f^{(1)}),\mathscr{S}(f^{(2)})\}=\{S_1,S_2,S_3\}.$
    Now $\zeta=\omega=e^{\frac{2\pi \mathfrak{i}}{3}}.$
    By \eqref{eq: cyclic group merging}, 
    \begin{equation}\label{eq: cyclic group n=3}
        f^{(0)}+\omega f^{(1)}+\omega^2 f^{(2)}=(1+\omega+\omega^2)g+(1+\omega^2+\omega^4)h=0.
    \end{equation}
    Notice that $y=0011$ is in $S_1$ but not in $S_2\cup S_3.$
    Evaluate \eqref{eq: cyclic group n=3} at $y=0011,$ LHS is nonzero, contradiction.
\end{proof}

\begin{lemma}\label{lm: cyclic group support size 2}
    Fix an arbitrary pair of indices $\{i,j\}$, $|\mathscr{S}(f_{ij}^{00})|=|\mathscr{S}(f_{ij}^{11})|=2.$ 
\end{lemma}
\begin{proof}
    Follow the notations in the proof of \cref{lm: common matching Cn},
    let $f^{(k)}=\partial_{ij}^{b^k}f,$ $g=f_{ij}^{00}$ and $h=f_{ij}^{11}.$
    Without loss of generality, we assume $\{i,j\}=\{5,6\}$.
    By \cref{lm: common matching Cn}, we may assume that for every $k\in[n]$, $f^{(k)}(x_1,x_2,x_3,x_4)=\lambda_k b^{^{k_1}}(x_1,x_2)b^{k_2}(x_3,x_4)$ for some $\lambda_k\neq 0$ and $k_1,k_2\in [n]$.
    Define $\widetilde{f^{(k)}}(x,y)=f^{(k)}(x,x,y,y),\,\widetilde{g}(x,y)=g(x,x,y,y),\,\widetilde{h}(x,y)=h(x,x,y,y)$ for every $x,y\in \{0,1\}.$
    By \eqref{eq: cyclic group merging}, we have $\widetilde{f^{(k)}}=\widetilde{g}+\zeta^k \widetilde{h}.$
    Notice that $$M(\widetilde{f^{(k)}})=\lambda_k\left[\begin{matrix}
        b^{(k_1)}(0,0)\\
        b^{(k_1)}(1,1)
    \end{matrix}\right]
    \left[\begin{matrix}
        b^{(k_2)}(0,0)&
        b^{(k_2)}(1,1)
    \end{matrix}\right].$$
    So $\mathrm{rank}(M(\widetilde{g}+\zeta^k\widetilde{h}))=\mathrm{rank}(M(\widetilde{f^{(k)}}))=1.$
    Hence $\det(M(\widetilde{g}+\zeta^k\widetilde{h}))=0.$
    Notice that $\det(M(\widetilde{g}+t\widetilde{h}))$ is a polynomial in $t$ of degree at most $2$, and it vanishes at $\zeta^k$, for all $k\in [n]$.
    Since $n\ge 3$, we have $\det(M(\widetilde{g}+t\widetilde{h}))\equiv 0.$
    \begin{claim}\label{claim: cyclic group rank 1}
        $\mathrm{rank}(M(\widetilde{g}))=\mathrm{rank}(M(\widetilde{h}))=1.$
    \end{claim}
    \begin{claimproof}{\cref{claim: cyclic group rank 1}}
        Let $t=0$ in $\det(M(\widetilde{g}+t\widetilde{h}))= 0$, we have $\det(M(\widetilde{g}))=0.$
        Hence $\mathrm{rank}(M(\widetilde{g}))\le 1.$
        Let $t\to +\infty$ in $\det(M(\widetilde{g}+t\widetilde{h}))= 0$, we have $\det(M(\widetilde{h}))=0$ and $\mathrm{rank}(M(\widetilde{h}))\le 1.$
        Next we show $\widetilde{g}\neq 0$ and $\widetilde{h}\ne 0.$
        By the assumption $f^{(k)}=\lambda_k b^{^{k_1}}(x_1,x_2)b^{k_2}(x_3,x_4)$, we have $\mathscr{S}(f^{(k)})=S$ for every $k\in [n]$, where $S=\{0000,0011,1100,1111\}.$
        By \eqref{eq: cyclic group merging}, we have 
        $$g=\frac{\zeta^{k_2}f^{(k_1)}-\zeta^{k_1}f^{(k_2)}}{\zeta^{k_2}-\zeta^{k_1}},\quad \text{for every distinct }k_1,k_2\in[n].$$
        So $\mathscr{S}(g)\subseteq S.$
        By assumption \eqref{eq: arity 6 assumption}, $f$ satisfies {\sc 2nd-Orth}.
        So $g=f_{ij}^{00}\neq 0.$
        It follows that $\widetilde{g}\neq 0$ by the definition of $\widetilde{g}$ and $\mathscr{S}(g)\subseteq S.$
        Thus, $\mathrm{rank}(M(\widetilde{g}))=1.$
        Similarly, we can show $\widetilde{h}\neq 0$ and $\mathrm{rank}(M(\widetilde{h}))=1.$
    \end{claimproof}
    By \cref{claim: cyclic group rank 1}, we may write $M(\widetilde{g})=uv^{\tt T}$ and $M(\widetilde{h})=wz^{\tt T}$, where $u,v,w,z\in \mathbb{C}^2.$
    If $u$ and $w$, $v$ and $z$ are both linearly independent, then $$M(\widetilde{g}+\widetilde{h})=\left[\begin{matrix}
        u & w
    \end{matrix}\right]
    \left[\begin{matrix}
        v^{\tt T} \\ z^{\tt T}
    \end{matrix}\right]$$
    has rank two, contradicting $\det(M(\widetilde{g}+t\widetilde{h}))\equiv 0.$
    Thus we may assume $u=w$ or $v=z$.
    After renaming variables, we may always assume $u=w.$
    Since $f$ satisfies {\sc 2nd-Orth}, we have $|g|=|h|$ and $\braket{g,h}=0.$
    So
    $$
    \braket{g,h}=(|u_0|^2+|u_1|^2)(v_0\overline{z_0}+v_1\overline{z_1})=0.
    $$
    Clearly $u\neq 0$, otherwise $\widetilde{g}=0.$
    So $v_0\overline{z_0}+v_1\overline{z_1}=0.$
    Also, $|g|=|h|$ gives $|v_0|^2+|v_1|^2=|z_0|^2+|z_1|^2.$ 
    \begin{claim}\label{claim: cyclic group equal norm}
        $|v_0|=|z_1|, |v_1|=|z_0|.$
    \end{claim}
    \begin{claimproof}{\cref{claim: cyclic group equal norm}}
        Let $r=(-\overline{v_1},\overline{v_0})^{\tt T}$, then $\braket{v,r}=0.$
        Since $\braket{v,z}=0$ and $v\neq 0$, there exists some $\lambda\in \mathbb{C}$ such that $z=\lambda r.$
        Substitute into $|v_0|^2+|v_1|^2=|z_0|^2+|z_1|^2$, we obtain $|v_0|^2+|v_1|^2=|\lambda|^2(|r_0|^2+|r_1|^2)=|\lambda|^2(|v_1|^2+|v_0|^2)$.
        So $|\lambda|=1.$
        Substituting back, we have $|v_0|=|z_1|$ and $|v_1|=|z_0|$.
    \end{claimproof}
    Now $\widetilde{f^{(k)}}=\widetilde{g}+\zeta^k\widetilde{h}=u(v+\zeta^k z)^{\tt T}$.
    So $f^{(k)}(x_1,x_2,x_3,x_4)=u(x_1,x_2)u'(x_3,x_4)$, where $u(0,0)=u_0$, $u(1,1)=u_1$, $u(0,1)=u(1,0)=0$; $u'(0,0)=v_0+\zeta^k z_0$, $u'(1,1)=v_1+\zeta^k z_1$ and $u'(0,1)=u'(1,0)=0.$
    By assumption \eqref{eq: arity 6 assumption}, both $u$ and $u'$ are proportional to some binary signatures in $B(f)$.
    Hence $|u_0|=|u_1|\neq 0$, and $|v_0+\zeta^k z_0|=|v_1+\zeta^k z_1|,\,\forall k\in [n].$
    Square both sides, we have
    \begin{equation*}
        |v_0|^2+2\Re(\zeta^k\overline{v_0}z_0)+|z_0|^2=|v_1|^2+2\Re(\zeta^k\overline{v_1}z_1)+|z_1|^2,\quad \forall k\in[n].
    \end{equation*}
    Using \cref{claim: cyclic group equal norm}, we have 
    \begin{equation*}
        \Re(\zeta^k(\overline{v_0}z_0-\overline{v_1}z_1))=0,\quad \forall k\in [n].
    \end{equation*}
    So $\overline{v_0}z_0-\overline{v_1}z_1=0.$
    Combining $v_0\overline{z_0}+v_1\overline{z_1}=0$, we have $\overline{v_0}z_0=\overline{v_1}z_1=0.$
    Since both $v$ and $z$ are non-zero, after possibly exchanging the coordinates, $v=(v_0,0)$ and $z=(0,z_1)$ with $v_0z_1\neq 0.$
    Consequently, $\widetilde{g}=uv^{\tt T}$ has two non-zero entries.
    By the definition of $\widetilde{g}$, we have $\mathscr{S}(g)=\mathscr{S}(\widetilde{g})=2$.
    Similarly, $\mathscr{S}(h)=2.$
    The lemma is proved.
\end{proof}

\begin{lemma}\label{lm: cyclic group even support}
    For any $x\in\mathscr{S}(f)$, $\mathrm{wt}(x)$ is even.
\end{lemma}
\begin{proof}
    Suppose $a=(a_1,a_2,a_3,a_4,a_5,a_6)\in \mathscr{S}(f)$ but $\mathrm{wt}(a)$ is odd.
    After renaming the variables, we may assume $a_1=a_2,a_3=a_4,a_5\neq a_6.$
    Consider merging $x_1$ and $x_2$ using $b^s\in B(f)$, and merging $x_3$ and $x_4$ using $b^t\in B(f)$.
    The resulting signature $\partial_{12}^{b^s}\partial^{b^t}_{34}f$ is also in $B(f)$ by assumption \eqref{eq: arity 6 assumption}.
    So $\partial_{12}^{b^s}\partial^{b^t}_{34}f(a_5,a_6)=0.$
    Thus,
    \begin{equation}\label{eq: cyclic group even support}
        \sum_{i,j\in\{0,1\}}(\zeta^s)^i(\zeta^t)^jf(i,i,j,j,a_5,a_6)=\partial_{12}^{b^s}\partial^{b^t}_{34}f(a_5,a_6)=0.
    \end{equation}
    Pick $s_1,s_2,t_1,t_2\in [n]$ such that $s_1\neq s_2$ and $t_1\neq t_2.$
    Then \eqref{eq: cyclic group even support} gives that 
    \begin{equation*}
        \left[\begin{matrix}
            1 & \zeta^{t_1} & \zeta^{s_1} & \zeta^{s_1t_1}\\
            1 & \zeta^{t_2} & \zeta^{s_1} & \zeta^{s_1t_2}\\
            1 & \zeta^{t_1} & \zeta^{s_2} & \zeta^{s_2t_1}\\
            1 & \zeta^{t_2} & \zeta^{s_2} & \zeta^{s_2t_2}\\
        \end{matrix}\right]
        \left[\begin{matrix}
            f(0,0,0,0,a_5,a_6)\\
            f(0,0,1,1,a_5,a_6)\\
            f(1,1,0,0,a_5,a_6)\\
            f(1,1,1,1,a_5,a_6)\\
        \end{matrix}\right]
        =\left[\begin{matrix}
            0\\
            0\\
            0\\
            0\\
        \end{matrix}\right].
    \end{equation*}
    The coefficient matrix is 
    $\left[\begin{smallmatrix}
        1 & \zeta^{s_1}\\
        1 & \zeta^{s_2}
    \end{smallmatrix}\right]\otimes\left[\begin{smallmatrix}
        1 & \zeta^{t_1}\\
        1 & \zeta^{t_2}
    \end{smallmatrix}\right]$, which has rank 4.
    Thus $f(a)=0$, contradiction.
\end{proof}

\begin{lemma}\label{lm: cyclic group impossible}
    Under assumption \eqref{eq: arity 6 assumption}, the standard case of $B(f)\cong C_n (n\ge 3)$ cannot occur.
\end{lemma}
\begin{proof}
    We count the number of pairs \((x,\{i,j\})\) such that $f(x)\ne0\text{ and }x_i=x_j.$
    There are fifteen choices of \(\{i,j\}\).
    For each pair, exactly four supported words satisfy $x_i=x_j$ by \cref{lm: cyclic group support size 2}.
    Thus the total number of incidences is
    $15\times4=60.$

    On the other hand,
    let \(e=\#\{x\in \mathscr{S}(f)\mid \mathrm{wt}(x)=0  \text{ or } 6\}\)
    and
    \(m=\#\{x\in \mathscr{S}(f)\mid \mathrm{wt}(x)=2\}\). 
    Clearly
    $0\le e\le2.$
    A word of weight \(0\) or \(6\) has $\binom{6}{2}=15$ equal pairs $\{x_i,x_j\}$.
    A word of weight \(2\) has $\binom22+\binom42=1+6=7$ equal pairs. 
    A word of weight \(4\) has the same number.
    By \cref{lm: cyclic group even support}, the total number of pairs $(x,\{i,j\})$ such that $f(x)\neq 0$ and $x_i=x_j$ is
    $15e+7m.$
    So $60=15e+7m$,
    $4\equiv e\pmod7.$
    Contradicting $0\le e\le 2.$
\end{proof}

\subsubsection{Non-standard Case}
The non-standard case can be handled similarly to the standard case by the following lemma.
We only sketch the proof and focus on differences from the standard case.
\begin{lemma} 
    Under assumption \eqref{eq: arity 6 assumption}, the non-standard case of $B(f)\cong C_n (n\ge 3)$ cannot occur.
\end{lemma}
\begin{proof}
    Fix a pair of indices $\{x_i,x_j\}.$
    Let $f^{(k)}=\partial_{ij}^{b^k}f,\, g=f_{ij}^{01},\, h=f_{ij}^{10}.$
    Without loss of generality, we assume $\{i,j\}=\{5,6\}.$
    We have 
    $
        f^{(k)}=g+\zeta^k h.
    $
    For every $k\in[n]$, $\mathscr{S}(f^{(k)})$ is one of the following:
    \begin{equation*}
        \begin{aligned}
        S_1&=\{0101,0110,1001,1010\}, \text{ or}\\
        S_2&=\{0011,0110,1001,1100\},\text{ or}\\
        S_3&=\{0011,0101,1010,1100\}.
        \end{aligned}
    \end{equation*}
    Again, if there exists two distinct $k,k'\in[n]$ such that $\mathscr{S}(f^{(k)})=\mathscr{S}(f^{(k')})=S$, then $\mathscr{S}(f^{(t)})=S$ for any $t\in[n]$.
    For $n\ge 4$, this is true by the pigeon hole principle.
    For $n=3$,
    unlike in \cref{lm: common matching Cn}, we require an exceptional argument here.
    Suppose $\{\mathscr{S}(f^{(0)}),\mathscr{S}(f^{(1)}),\mathscr{S}(f^{(2)})\}=\{S_1,S_2,S_3\}.$
    Then 
    \[\{f^{(0)},f^{(1)},f^{(2)}\}=\{\lambda_0b^{r_1}(x_1,x_2)b^{r_2}(x_3,x_4),\lambda_1b^{s_1}(x_1,x_3)b^{s_2}(x_2,x_4),\lambda_2b^{t_1}(x_1,x_4)b^{t_2}(x_2,x_3)\}\]
    for some $\lambda_0,\lambda_1,\lambda_2\neq 0$ and $r_1,r_2,s_1,s_2,t_1,t_2\in\{0,1,2\}.$
    By
    $
        f^{(k)}=g+\omega^k h
    $
    where $\omega=e^{\frac{2\pi \mathfrak{i}}{3}}$, we have
    $f^{(0)}+\omega f^{(1)}+\omega^2 f^{(2)}=0.$
    Absorbing the coefficients $\lambda_k$, we have
    $$
    \alpha b^{r_1}(x_1,x_2)b^{r_2}(x_3,x_4)+\beta b^{s_1}(x_1,x_3)b^{s_2}(x_2,x_4)+\gamma b^{t_1}(x_1,x_4)b^{t_2}(x_2,x_3)=0
    $$
    for some $\alpha,\beta,\gamma\neq0.$
    Evaluate at $(x_1,x_2,x_3,x_4)=0110,0011$ and $0101$ respectively, we have
    \begin{equation*}
        \begin{aligned}
            &\alpha \omega^{r_2}+\beta \omega^{s_2}=0,\\
            & \beta+\gamma=0,\\
            &\alpha+\gamma \omega^{t_2}=0.
        \end{aligned}
    \end{equation*}
    The second equation gives $\beta=-\gamma$.
    Plug into the first, $\alpha=\gamma\omega^{s_2-r_2}.$
    Plug into the third, $\omega^{t_2}=-\omega^{s_2-r_2}.$
    Cube both sides, we have $-1=1$, contradiction.
    Thus, the same statement (after replacing $B(f)$ by $\widehat{B}(\widehat{f})$) in \cref{lm: common matching Cn} can be proved in the non-standard case.
    Similar to \cref{lm: cyclic group support size 2}, we can prove $|\mathscr{S}(f_{ij}^{01})|=|\mathscr{S}(f_{ij}^{10})|=2$ for any pair of indices $\{i,j\}.$

    Instead of proving $f$ is supported on even Hamming weights like \cref{lm: cyclic group even support}, we will show that for any $x\in \mathscr{S}(f)$, $\mathrm{wt}(x)=0,6$ or $3$.
    Pick a non-constant $x\in \mathscr{S}(f)$.
    Choose two positions \(i,j\) on which \(x_i\ne x_j\), and let \(y\) denote the remaining four bits.
    Since $f(x)\neq 0$, the degree one polynomial
    \(
    P_y(t)=f_{ij}^{01}(y)+tf_{ij}^{10}(y)
    \)
    is not identically zero.
    Therefore, there exists some $k\in[n]$ such that $f^{(k)}(y)=P_y(\zeta^k)\neq 0.$
    Since $\mathscr{S}(f^{(k)})=S_1,S_2$ or $S_3$, we have $\mathrm{wt}(y)=2.$
    Hence $\mathrm{wt}(x)=3.$
    We have proved $\mathscr{S}(f)\subseteq \{000000,111111\}\cup \{x\mid x\in \{0,1\}^6,\mathrm{wt}(x)=3\}$.

    Now we count the number of pairs $(x, \{i,j\})$ such that $f(x)\neq 0$ and $x_i\ne x_j$.
    Again there are fifteen choices of \(\{i,j\}\), and
    for each pair, exactly four supported words satisfy $x_i=x_j$.
    Thus the total number of incidences is
    $15\times4=60.$
    On the other hand, the two constant words contribute zero unequal pairs, and
    a weight-three word has $3\times 3=9$ unequal pairs.
    So $60=9m$, where $m\in \mathbb{Z}$ is the number of $x\in \mathscr{S}(f)$ such that $\mathrm{wt}(x)=3.$
    This is a contradiction.
\end{proof}


